\documentclass[11pt,a4paper]{article}
\usepackage[margin=1in]{geometry}
\usepackage{amsmath,amssymb}
\usepackage{parskip}
\usepackage{hyperref}

\title{\textbf{Pre-Experiment Overview:\\ Electron Diffraction}}
\author{IBRAHIM H.I. ABUSHAWISH \\ Istanbul University, Department of Physics}
\date{}

\begin{document}
\maketitle

\section*{Purpose of the Experiment}
According to the "ATOMIC AND MOLECULAR PHYSICS LABORATORY-2nd.pdf" manual, the primary goal of this lab is to direct free electrons through a polycrystalline graphite layer to observe their diffraction. By doing this, we are trying to physically prove that electrons possess wave characteristics. On a practical level, we also use the resulting diffraction patterns to calculate the specific lattice constants, $d_1$ and $d_2$, of the graphite crystal.

\section*{Importance and Significance}
This experiment is highly significant because it provides a direct, hands-on confirmation of Louis de Broglie's 1923 hypothesis, which proposed that massive particles inherently have a wave nature. Historically, this exact concept was confirmed by Davisson and Germer, eventually earning de Broglie a Nobel Prize.

In the lab setup, as the accelerated electrons hit the hexagonal structure of the graphite, they scatter off different atomic planes. This creates an interference pattern visible as two bright concentric rings on a fluorescent screen, a phenomenon known as Bragg diffraction. The most crucial part of the experiment---the part that actually proves the underlying physics---is how these rings react to voltage changes. When the accelerating anode voltage is decreased, the electrons lose momentum, which causes their de Broglie wavelength to increase. This increased wavelength directly results in a larger angular dispersion, meaning the diffraction rings visibly get wider. Seeing the ring sizes change in perfect mathematical agreement with our voltage adjustments is the ultimate, visible proof of wave-particle duality.

\end{document}
